\section{The Shared Test Files and the Protocol}
\label{sec:benchmark}

\paragraph{What the field evaluates on.}
\label{sec:substrate}
The EarningsCall (EC) benchmark of \citet{qin2019what} and the two MAEC cohorts
are scored on split and label files that VolTAGE and KeFVP redistribute. We
audited those files (Table~\ref{tab:substrate}). The VolTAGE and KeFVP test
files are byte-identical, so the two releases reuse the same label file. Five
of the six boundaries between training, validation and test have a zero-day gap (the sixth has three days)
while the targets reach 30 days past each call. On the time-ordered splits the field reports, 80.4\% (EC), 44.2\% (MAEC-15)
and 55.7\% (MAEC-16) of test calls are from companies that also appear in
training. Any time-ordered split allows that, and no paper reports what its
model does when the company is held out. The original paper defines its horizons in calendar
days, but most later work describes them as trading days (21 sessions rather than
30 at $\tau{=}30$). Finally, the shipped MAEC labels correlate only 0.52--0.87 at $\tau\ge7$
(0.25--0.66 at $\tau{=}3$) with calendar-day targets we recompute from prices
under the same call-date anchor, and no single stated rule reproduces them.

% Curated from data/results/result_table_5r.csv and result_table_5r_labels.csv
% (runs 20261006T202254Z / 20261006T202821Z-audit-substrate-48b9432a). Overlap = test calls whose
% ticker appears in train; embargo = calendar days between the last train (val)
% date and the first val (test) date. Ours = committed splits, same metrics.
\begin{table}[t]
\centering
\small
\begin{tabular}{lrrrrr}
\toprule
Benchmark / split & calls (tr/va/te) & overlap & embargo tr$\to$va & va$\to$te & label defect \\
\midrule
EC (Qin \& Yang)         & 392 / 56 / 112  & 80.4\% & 0 d & 0 d & 137 / 16{,}800 zero cells \\
MAEC-15 (KeFVP files)    & 535 / 76 / 154  & 44.2\% & 0 d & 0 d & 4 binary cols at $\tau$ \\
MAEC-16 (KeFVP files)    & 980 / 140 / 280 & 55.7\% & 0 d & 3 d & 4 binary cols at $\tau$ \\
\midrule
ours: FinCall temporal        & 1489 / 228 / 483 & 91.3\% & 47 d & 47 d & --- \\
ours: FinCall ticker-disjoint & 1706 / 243 / 484 & \textbf{0.0\%} & n/a & n/a & --- \\
ours: MAEC temporal           & 1666 / 69 / 529 & 79.4\% & 101 d & 48 d & --- \\
ours: MAEC ticker-disjoint    & 1884 / 271 / 533 & \textbf{0.0\%} & n/a & n/a & --- \\
ours: Earnings25 ticker-disjoint & 333 / 47 / 95 & \textbf{0.0\%} & n/a & n/a & --- \\
\bottomrule
\end{tabular}
\caption{\textbf{Substrate audit of the shared benchmark files} (Result Table 5R,
regenerated by \texttt{ecvol audit substrate} from SHA-256-manifested copies of the
VolTAGE and KeFVP repositories). The VolTAGE and KeFVP split files are
byte-identical. Overlap is a property every temporal split permits, ours included;
the point is that no paper in the lineage reports the disjoint condition. The
zero-day embargos leak 30-day target windows across every boundary. Shipped
MAEC labels join our calendar-day, call-date-anchored targets on 77--79\% of test calls at
$\tau\ge7$ and correlate 0.52--0.87 with them (0.25--0.66 at $\tau{=}3$).}
\label{tab:substrate}
\end{table}


\paragraph{Corpora and company identity.}
\textbf{FinCall-Surprise} \citep{fincall2025} is the primary corpus, with 2{,}688
calls from 2019--2021 and MP3 audio, transcripts and slides, but no ticker, date,
CIK or speaker names. We recover the company and the date from the transcript
text and the slide metadata, check them against price data, and release the
table with its 80 manual corrections. 2{,}496 calls resolve (92.9\%, 95.3\% of
the earnings-type calls), and a 60-call manual audit passes 60/60. Nine resolved
calls (0.4\%) carry a pre-2019 date that their transcripts contradict, and all sit
in training segments and are flagged in the release. \textbf{MAEC}
\citep{li2020maec} is the secondary corpus, with 3{,}443 calls from 2015--2018 and sentence-aligned audio features. The release we
could access ships no raw audio.
\textbf{Earnings25} \citep{jiang2026earnings25} has 514 S\&P~500 calls in
the release we ingested, 501 of them from October to December 2025, with webcast
MP3 at 64\,kbps for 281 calls and at most 24\,kbps for 232. 476 of its calls map
to a ticker and join to prices, and 475 carry a target.

\paragraph{Targets.}
Daily adjusted prices come from Yahoo Finance through the yfinance package,
whose terms do not permit redistribution, and were checked against
Tiingo on 50 tickers. We release derived volatility targets, never the price
series. Let $r_t$ be the simple return of the adjusted close on session $t$, counted from day~0, the last
session before the call's information is public. The main target is
$v_{\mathrm{post}}(\tau)=\ln\sqrt{\frac{1}{n}\sum_{t=1}^{n}(r_t-\bar r)^2}$ for
$\tau\in\{3,7,15,30\}$, and $v_{\mathrm{pre}}(\tau)$ is the same quantity over the
sessions ending at day~0. Under the \textbf{trading-day} definition $n=\tau$. Under
the \textbf{calendar-day} definition of \citet{qin2019what} the window holds every
session within $\tau$ calendar days, so $n\approx0.7\tau$. We drop any window
with fewer than two sessions and record the reason, which removes 42.5\% of
FinCall's calls at $\tau{=}3$ calendar days. $v_{\mathrm{post}}$ correlates only
0.78 between the definitions at $\tau{=}3$ and the persistence MSE is 44--63\%
higher under calendar days at $\tau{=}3$ (31--41\% at longer horizons), so a published MSE means little unless its
definition is stated (Appendix Table~\ref{tab:convention}). Every result below
uses trading days, and the definition is a column in every result table. We add two targets
that are harder to game with company identity. The \textbf{change}
$\Delta v=v_{\mathrm{post}}-v_{\mathrm{pre}}$ subtracts the company's own level,
and the \textbf{HAR-residual} subtracts a HAR-RV forecast. After joining to
prices, FinCall has 2{,}433 calls with a target (97.5\%) and MAEC 2{,}688 of 3{,}443
(78.1\%, the rest lacking usable price history).

\paragraph{When the information became public.}
\label{sec:timing}
\citet{scss2024} showed that most calls in EC and MAEC are released before the
open and that the shipped labels ignore this, which their DEC corpus corrects
with a vendor flag. Almost no transcript states a clock time, so we use the
EDGAR acceptance time of the earliest 8-K filing with Item~2.02 (the earnings
release) within three days before to one day after the call as a proxy for the
release time. That covers 89.3\% of FinCall, 68.6\% of MAEC and 99.8\% of Earnings25. The
rest fall back, in recorded order, to Earnings25's call-start time, the
before/after-market flag of \citet{scss2024}, and an after-hours assumption. A
release before the open makes the call date itself day~1. Compared with the
call-date anchor that every released label set uses, \textbf{day~0 moves for 53\%
of FinCall, 56\% of MAEC and 72\% of Earnings25 calls}, almost always one session
earlier. Under the published anchor, the session in which the market reacts
therefore sits in the \emph{pre} window for most calls, and $\Delta v$ gains a
positive mean on the measured anchor ($+0.51$ at $\tau{=}3$ on FinCall). The
call-date targets ship as a labeled variant. On
Earnings25, which records call start times, the filing precedes the call for
98.9\% of calls (median 100 minutes). A press release can precede its 8-K.
Filings accepted within an hour after the open, where an earlier release could
move day~0, are 1.1\% of FinCall, 3.1\% of MAEC and 0.8\% of Earnings25 calls,
a sensitivity diagnostic rather than an upper bound on the disagreement. The
call can still be running during session~1, so these targets describe a release-aligned event window. A
forward-looking variant starts the post window one session later, so that no
target return begins before the call has ended (true for all 476 Earnings25
calls with a recorded start time).

\paragraph{Splits and release.}
\emph{Temporal} splits (about 70/10/20 in date order) leave a 30-trading-day gap
between segments so that no target window crosses a boundary.
\emph{Ticker-disjoint} splits keep every validation and test company out of
training. Their overlap is 0.0\%, against 91.3\% (FinCall) and 79.4\% (MAEC) on
our own temporal splits. Because they are not in date order, a
\emph{date-ordered company-disjoint} split applies both rules at once, at the
cost of small test and validation sets (97 and 22 FinCall calls). Appendix Table~\ref{tab:splits} gives calls,
companies, overlaps and dates for every split. The harness fails any feature that reads price data after
the call's as-of time, and leakage assertions run in CI. The release contains all
of the above, the derived-feature archives, the harness, and every result CSV with
its run record (resolved config, git commit, seeds and file hash).
